\documentclass[12pt]{article}

\usepackage{amsmath,amssymb}
\usepackage[margin=1in]{geometry}
\usepackage{hyperref}
\usepackage{booktabs}
\usepackage{microtype}

% ── notation shortcuts ─────────────────────────────────────────────────────────
\newcommand{\tauv}{\tau_{V}}
\newcommand{\vdisp}{\sigma_{\star}}
\newcommand{\klam}{k(\lambda)}
\newcommand{\Tigm}{T_{\mathrm{IGM}}}
\newcommand{\DL}{D_L}
\newcommand{\Msun}{M_{\odot}}
\newcommand{\Lsun}{L_{\odot}}
\newcommand{\flam}{f_\lambda}
\newcommand{\Flam}{F_\lambda}
\newcommand{\Mstar}{M_{\star}}
\newcommand{\R}{\mathbf{R}}
\newcommand{\cc}{\mathbf{c}}
\newcommand{\C}{c}                 % speed of light

\title{FastSpecFit Continuum Model}
\author{John Moustakas%
  \thanks{Continuum-model derivations and implementation assisted by
    Claude Code (Claude Sonnet 4.6, Anthropic).}}
\date{\today}


\begin{document}
\maketitle

\section{Introduction}

This note documents the stellar continuum forward model used by
\textsc{FastSpecFit} to fit DESI galaxy spectra and broadband photometry.
The pipeline operates in two modes:
\begin{itemize}
  \item \textbf{\texttt{fastphot}} — photometry-only mode: fits broadband
    spectral energy distributions (SEDs) to Legacy Survey DR9/DR10
    photometry with no spectroscopic data.
  \item \textbf{\texttt{fastspec}} — full spectrophotometric mode: jointly
    fits spectra from all three DESI cameras (B, R, Z) and broadband
    photometry simultaneously.
\end{itemize}

The topics covered are: (i)~the stellar population synthesis (SPS) template
library and how it is constructed; (ii)~the dust attenuation and emission
models; (iii)~the redshift, IGM, and flux-scaling factors applied to the
templates; (iv)~velocity-dispersion broadening; (v)~photometric and
spectroscopic synthesis; (vi)~the fitting algorithms, including the variable
projection approach used in the primary fits; and (vii)~the nonparametric
smooth continuum correction, Monte Carlo uncertainty estimation, and physical
output quantities.


\section{Stellar Population Synthesis Templates}
\label{sec:templates}

\subsection{Template construction}
\label{ssec:build}

The SPS template library is built with the \texttt{python-fsps} interface
to FSPS (Flexible Stellar Population Synthesis; Conroy, Wechsler \&
van Dokkum 2009; Conroy \& Gunn 2010) using:
\begin{itemize}
  \item \textbf{IMF}: Chabrier (2003).
  \item \textbf{Isochrones}: MIST (Choi et al.\ 2016).
  \item \textbf{Stellar library}: C3K\_a (Conroy et al.), covering
    2750--9100\,\AA\ at resolution $R \equiv \lambda/\mathrm{FWHM} = 3000$
    (equivalent to ${\sim}42$\,km\,s$^{-1}$).
  \item \textbf{Star-formation history (SFH)}: constant SFR within each of
    five variable-width lookback-time bins spanning 30\,Myr to 13.7\,Gyr.
    The bin boundaries follow the ``continuity'' prescription of Leja et al.\
    (2017): the first two bins have fixed widths of $\Delta\log_{10}(t/\mathrm{yr})
    = 7.48$ and $8.0$, and the remaining three are logarithmically spaced up
    to $0.85\,t_{\mathrm{univ}}$.
  \item \textbf{Metallicity}: one or more values of $\log_{10}(Z/Z_\odot)$,
    defaulting to solar.
  \item \textbf{Nebular emission}: included via the FSPS built-in
    photoionization grid; a separate set of templates without nebular emission
    is produced simultaneously.
\end{itemize}

Each template spectrum is normalized so that one solar mass of stars formed
at the observed age would produce the tabulated flux at 10\,pc:
\begin{equation}
  \Flam \;[\mathrm{erg\,s^{-1}\,cm^{-2}\,\text{\AA}^{-1}\,\Msun^{-1}}]
  \;=\; \frac{L_\odot}{4\pi (10\,\mathrm{pc})^2}\,\frac{L_\lambda}{L_\odot}.
\end{equation}
This normalization means that the fitted coefficients $\cc$ have units of
$\Msun$ (subject to additional scaling by the mass normalization factor
$\mathcal{M}_0 = 10^{10}\,\Msun$).

\subsection{Wavelength grid}
\label{ssec:grid}

Templates are resampled from the native FSPS grid to a hybrid wavelength
array.  In the range 2750--9100\,\AA\ the pixel size is constant in velocity,
$\Delta v = 25$\,km\,s$^{-1}$ (i.e., a uniform $\Delta\ln\lambda$); outside
this range the native FSPS grid is retained.  The constant velocity pixel
enables Gaussian convolution via a fixed-width kernel in log-$\lambda$ space
(Section~\ref{sec:vdisp}).

\subsection{Auxiliary templates}
\label{ssec:aux}

In addition to the stellar continua, the template file bundles:
\begin{itemize}
  \item \textbf{Dust emission} — Draine \& Li (2007; DL07) interstellar
    dust model, interpolated over $q_\mathrm{PAH}$, $U_\mathrm{min}$, and
    $\gamma$ parameters and normalized to unit bolometric luminosity
    (Section~\ref{ssec:dustem}).
  \item \textbf{AGN torus} — Nenkova et al.\ (2008) clumpy-torus model,
    indexed by optical depth $\tau$; used optionally for AGN host fitting.
  \item \textbf{Fe emission} — Vestergaard \& Wilkes (2001) empirical
    UV/optical iron emission template for broad-line AGN fitting.
\end{itemize}


\section{Dust Attenuation and Emission}
\label{sec:dust}

\subsection{Stellar attenuation}
\label{ssec:atten}

Stellar continuum attenuation is modeled as a uniform foreground screen
using a power-law attenuation curve,
\begin{equation}
  \klam \;=\; \left(\frac{\lambda}{5500\,\text{\AA}}\right)^{-0.7},
  \label{eq:klam}
\end{equation}
with a single free parameter, the $V$-band optical depth $\tauv$.  The
attenuated model spectrum is
\begin{equation}
  \Flam^{\mathrm{att}}(\lambda)
  \;=\; \Flam(\lambda)\,\exp\!\bigl(-\tauv\,\klam\bigr).
  \label{eq:atten}
\end{equation}
The power-law index $-0.7$ is the Charlot \& Fall (2000) birth-cloud slope
and is also consistent with the Calzetti et al.\ (2000) attenuation law
shortward of $\sim6000$\,\AA.  The single-parameter screen is intentionally
simple; age-dependent attenuation is a planned future extension.

\subsection{Dust emission and energy balance}
\label{ssec:dustem}

Absorbed photons are re-emitted in the infrared.  The DL07 template
$\Flam^{\mathrm{dust}}(\lambda)$ (normalized to unit bolometric luminosity)
is scaled by the absorbed bolometric luminosity,
\begin{equation}
  L_{\mathrm{abs}}
  \;=\; \int_0^\infty \Flam(\lambda)
        \bigl[1 - \exp(-\tauv\,\klam)\bigr]\,d\lambda,
  \label{eq:labs}
\end{equation}
so that the total (attenuated stellar $+$ re-emitted dust) model conserves
bolometric luminosity:
\begin{equation}
  \boxed{
  \Flam^{\mathrm{tot}}(\lambda)
  \;=\; \Flam^{\mathrm{att}}(\lambda)
        \;+\; L_{\mathrm{abs}}\,\Flam^{\mathrm{dust}}(\lambda).
  }
  \label{eq:etot}
\end{equation}
Dust self-absorption is neglected.  The DL07 template is built with
default parameters $q_\mathrm{PAH} = 3.5$, $U_\mathrm{min} = 1.0$, and
$\gamma = 0.01$ (fraction of dust heated by photodissociation regions).


\section{Redshift, IGM, and Flux Scaling}
\label{sec:zfactors}

For an object at redshift $z$ with luminosity distance $\DL$ (Mpc), the
observed flux density is related to the rest-frame template by
\begin{equation}
  \flam^{\mathrm{obs}}(\lambda)
  \;=\;
  \Tigm(\lambda,z)\,
  \frac{\mathcal{F}_0\,\mathcal{M}_0}{4\pi (3.086\times10^{24}\,\DL)^2}
  \,\frac{\Flam^{\mathrm{tot}}(\lambda/(1+z))}{1+z},
  \label{eq:zfactor}
\end{equation}
where $\Tigm(\lambda,z)$ is the wavelength-dependent IGM transmission from
the Inoue et al.\ (2014) model (Ly$\alpha$ forest and Lyman-limit systems),
$\mathcal{F}_0 = 10^{17}$ is the flux normalization factor
(converting erg\,s$^{-1}$\,cm$^{-2}$\,\AA$^{-1}$ to the internal units
$10^{-17}$\,erg\,s$^{-1}$\,cm$^{-2}$\,\AA$^{-1}$), and
$\mathcal{M}_0 = 10^{10}\,\Msun$ is the mass normalization.

In practice these factors are combined into a single array,
\begin{equation}
  Z_i \;=\; \Tigm(\lambda_i,z)\,
            \frac{\mathcal{F}_0\,\mathcal{M}_0}
                 {4\pi(3.086\times10^{24}\,\DL)^2\,(1+z)},
  \label{eq:zfactors}
\end{equation}
evaluated once at object initialization and cached as \texttt{zfactors}.
The redshifted template wavelength array $\lambda_i^z = \lambda_i^{\rm rest}(1+z)$
is similarly cached as \texttt{ztemplatewave}.


\section{Velocity Dispersion Broadening}
\label{sec:vdisp}

\subsection{Gaussian convolution in log-$\lambda$}
\label{ssec:conv}

On the constant-velocity wavelength grid (Section~\ref{ssec:grid}), a
Gaussian with line-of-sight velocity dispersion $\vdisp$\,(km\,s$^{-1}$)
corresponds to a kernel of fixed width
\begin{equation}
  \sigma_{\mathrm{pix}} \;=\; \vdisp / \Delta v,
  \label{eq:sigma_pix}
\end{equation}
where $\Delta v = 25$\,km\,s$^{-1}$ is the pixel size.  Convolution is
performed via FFT over the 2750--9100\,\AA\ subrange; outside this range the
native (coarser) IR pixels are retained unchanged.

\subsection{Pre-computed FFT cache}
\label{ssec:cache}

To avoid repeated FFT computation during fitting, the code pre-computes and
caches the Fourier transforms of all templates at velocities up to
$\vdisp^{\mathrm{max}} = 500$\,km\,s$^{-1}$.  Convolution at a requested
$\vdisp \leq \vdisp^{\mathrm{max}}$ then requires only a single inverse FFT
of the product of the cached template transform and the (analytic) Gaussian
transfer function.  For $\vdisp > \vdisp^{\mathrm{max}}$ the convolution is
performed directly in pixel space.

\subsection{Fitting the velocity dispersion}
\label{ssec:vdispfit}

The velocity dispersion is fitted only when the spectrum covers a sufficient
rest-frame wavelength range; specifically, when $\lambda_{\min}^{\mathrm{rest}}
< 3800$\,\AA\ and $\lambda_{\max}^{\mathrm{rest}} > 4800$\,\AA.  The fit
uses only the 3800--6000\,\AA\ rest-frame window, which captures the Ca\,H\&K
break, the $G$-band, Mg\,\textsc{b}, and Fe absorption features sensitive to
$\vdisp$.

Fitting proceeds in two stages:

\begin{enumerate}
  \item \textbf{Chi-squared grid scan.}  Templates (without nebular emission)
    are convolved at $N_{\mathrm{grid}} = 5$ evenly spaced values of $\vdisp$
    in the range $[75, 500]$\,km\,s$^{-1}$.  At each grid point a full
    continuum fit (Section~\ref{sec:fitting}) is performed and the chi-squared
    is recorded.  A grid result is accepted only if the peak-to-peak
    $\Delta\chi^2 \geq 25$ and the minimum does not fall at a boundary.
  \item \textbf{Maximum-likelihood refinement.}  If the grid scan succeeds,
    the full TRF optimizer (Section~\ref{ssec:trf}) is run with $\vdisp$ as
    an additional free parameter, initialized at the grid minimum.  If the
    scan fails, $\vdisp$ is fixed at the nominal value of 250\,km\,s$^{-1}$.
\end{enumerate}

The nominal dispersion of 250\,km\,s$^{-1}$ is also used whenever wavelength
coverage is insufficient; in that case the cached \texttt{flux\_nomvdisp}
templates (pre-convolved to 250\,km\,s$^{-1}$) are used directly without
additional convolution.


\section{Photometric and Spectroscopic Synthesis}
\label{sec:synthesis}

\subsection{Photometry}
\label{ssec:synthphot}

Synthetic AB photometry in filter $f$ is computed as
\begin{equation}
  m_f^{\mathrm{AB}} \;=\; -2.5\log_{10}\!\left(
    \frac{\int_0^\infty (\flam^{\mathrm{obs}} / h\nu)\,R_f(\lambda)\,d\lambda}
         {\int_0^\infty (f_{\nu,0} / h\nu)\,R_f(\lambda)\,d\lambda}
  \right),
  \label{eq:synthphot}
\end{equation}
where $R_f(\lambda)$ is the filter response curve (from
\texttt{speclite}) and $f_{\nu,0}$ is the AB zero-point flux density.
Numerically, the integrals are evaluated by trapezoidal quadrature over the
redshifted template grid; the per-filter response-weighted weights
(\texttt{phot\_batch\_weights}) are pre-computed once and reused for all
template evaluations during fitting.

The synthesized flux density $\flam$ in units of
$10^{-17}$\,erg\,s$^{-1}$\,cm$^{-2}$\,\AA$^{-1}$ is obtained from the
maggies $m$ via $\flam = m \times 3631 \times \nu^2/(\C \times 10^{17})$,
where $\nu$ is the filter effective frequency.

\subsection{Spectroscopy}
\label{ssec:synthspec}

The observed-frame model spectrum at DESI pixels is computed in two steps:
\begin{enumerate}
  \item \textbf{Resampling.}  The high-resolution template (on the native
    redshifted grid) is flux-conservatively resampled onto the per-camera
    DESI wavelength grid using trapezoidal rebinning (\texttt{trapz\_rebin}).
  \item \textbf{Resolution convolution.}  The resampled spectrum is multiplied
    by the DESI resolution matrix $\R$,
    \begin{equation}
      d_i^{\mathrm{model}} \;=\; \sum_j R_{ij}\,\tilde{f}_j,
    \end{equation}
    where $\tilde{f}_j$ is the resampled flux and $\R$ encodes the
    wavelength-dependent LSF in band format (${\sim}11$ diagonals).
    This step accounts for the instrumental line-spread function in the
    same way as the emission-line model (see the companion technote).
\end{enumerate}


\section{Fitting Algorithms}
\label{sec:fitting}

\subsection{Model parameterization}
\label{ssec:params}

The full continuum model for a galaxy is
\begin{equation}
  \flam^{\mathrm{model}}(\lambda;\,\tauv,\vdisp,\cc)
  \;=\; Z(\lambda)\,
        \Bigl[\cc \cdot \Phi(\lambda;\vdisp)\Bigr]^{\mathrm{att}}(\tauv),
  \label{eq:model}
\end{equation}
where $\Phi(\lambda;\vdisp)$ is the matrix of velocity-broadened template
spectra (shape $[N_{\mathrm{age}} \times N_{\mathrm{met}},\,N_\lambda]$),
$\cc \geq 0$ is the vector of non-negative template coefficients, and the
superscript ``att'' denotes the attenuation and dust-emission operations of
Section~\ref{sec:dust}.  The free parameters are $\tauv \in [0, 2]$,
$\vdisp \in [75, 500]$\,km\,s$^{-1}$ (if fitted), and $\cc \geq 0$.

The fitting objective is the weighted sum of squared spectroscopic and
photometric residuals,
\begin{equation}
  \chi^2(\tauv,\vdisp,\cc)
  \;=\;
  \sum_i \frac{\bigl(f_i^{\mathrm{obs}} - f_i^{\mathrm{model}}\bigr)^2}
              {\sigma_i^2}
  \;+\;
  \sum_k \frac{\bigl(F_{\lambda,k}^{\mathrm{obs}} - F_{\lambda,k}^{\mathrm{model}}\bigr)^2}
              {\sigma_{\!f,k}^2},
  \label{eq:chi2}
\end{equation}
where the first sum runs over unmasked spectral pixels and the second over
photometric bands with $\sigma_{\!f,k}^2 > 0$.

\subsection{Variable projection (VARPRO)}
\label{ssec:varpro}

When $\vdisp$ is fixed (either at its nominal value or after the vdisp fit),
the model~\eqref{eq:model} is \emph{linear} in $\cc$ for any fixed $\tauv$.
This separable structure is exploited via variable projection:

\begin{enumerate}
  \item \textbf{Outer loop}: a bounded scalar minimization over $\tauv$ using
    Brent's method (\texttt{scipy.optimize.minimize\_scalar}; typically
    10--15 function evaluations).
  \item \textbf{Inner solve}: for each trial $\tauv$, assemble the weighted
    design matrix
    \begin{equation}
      \Psi_{ik} \;=\; \sigma_i^{-1}\,[\mathbf{S}\,\Phi^{\mathrm{att}}]_{ik},
      \label{eq:Psi}
    \end{equation}
    where $\mathbf{S}$ represents either the spectroscopic synthesis
    (Section~\ref{ssec:synthspec}) or photometric synthesis
    (Section~\ref{ssec:synthphot}) operator (or both, stacked vertically),
    and solve the non-negative least-squares problem
    \begin{equation}
      \boxed{
      \cc^* \;=\; \arg\min_{\cc\,\geq\,0}
                  \|\Psi\,\cc - \mathbf{b}\|^2,
      }
      \label{eq:nnls}
    \end{equation}
    with $\mathbf{b}$ the vector of (weighted) observed data.
    This is solved with \texttt{scipy.optimize.nnls}.
\end{enumerate}

The VARPRO approach avoids finite-differencing over all $N_{\mathrm{age}}
+ N_{\mathrm{met}} + 1$ parameters and typically converges $10\times$ faster
than a direct TRF optimization over the full parameter vector.

\subsection{Trust-region reflective optimizer}
\label{ssec:trf}

When $\vdisp$ is a free parameter, the separability is broken and the full
parameter vector $(\tauv, \vdisp, \cc)$ is passed to
\texttt{scipy.optimize.least\_squares} with:
\begin{itemize}
  \item Method: \texttt{trf} (trust-region reflective).
  \item Jacobian scaling: \texttt{x\_scale='jac'} (Jacobian-column scaling)
    to improve conditioning across the disparate parameter scales.
  \item Solver: \texttt{tr\_solver='exact'} with regularization.
  \item Bounds: $\tauv \in [0,2]$, $\vdisp \in [75, 500]$\,km\,s$^{-1}$,
    $c_k \in [0, 10^6]$.
  \item Tolerances: $f_{\mathrm{tol}} = 10^{-6}$, $x_{\mathrm{tol}} = 10^{-10}$.
\end{itemize}

The Jacobian is computed by finite differencing.


\subsection{Aperture correction}
\label{ssec:apercorr}

DESI fibers subtend a finite sky area and typically capture only a fraction
of the total galaxy flux.  After the initial spectroscopic vdisp fit (using
the fiber spectrum), a per-band aperture correction is estimated as
\begin{equation}
  \eta_f \;=\; \frac{F_{\lambda,f}^{\mathrm{fiber,obs}}}
                     {F_{\lambda,f}^{\mathrm{model}}(\lambda_{\mathrm{eff},f})},
  \label{eq:apercorr}
\end{equation}
where $F_{\lambda,f}^{\mathrm{fiber,obs}}$ is the observed fiber photometry in
band $f$ and $F_{\lambda,f}^{\mathrm{model}}$ is the synthesized photometry from
the initial fit.  The median over bands, $\bar{\eta}$, is applied to the
spectrum ($\flam \to \bar{\eta}\,\flam$) and to the spectral uncertainties
($\sigma \to \sigma/\bar{\eta}$) before the full spectrophotometric fit.
This correction brings the spectral flux level into agreement with the total
broadband photometry.

\subsection{Full spectrophotometric fit}
\label{ssec:fullfit}

The final fit uses the VARPRO solver (Section~\ref{ssec:varpro}) with $\vdisp$
fixed at the value determined in Section~\ref{ssec:vdispfit}, using templates
that include stellar continuum $+$ nebular emission.  Spectroscopy
(aperture-corrected) and photometry are fitted simultaneously, so the
residual vector stacks spectral and photometric terms.  After convergence,
the best-fit model without nebular emission is built separately (with no
dust emission) to measure Dn(4000) and rest-frame continuum fluxes.


\section{Nonparametric Smooth Continuum}
\label{sec:smooth}

The parametric SED model \eqref{eq:model} may leave low-level systematic
residuals due to imperfect template matching, flux calibration errors, or
Galactic extinction uncertainties.  \textsc{FastSpecFit} estimates and
removes these residuals with a nonparametric smooth continuum computed from
the spectrum after subtracting the best-fit line-free SED model.

The algorithm operates per camera:
\begin{enumerate}
  \item Mask pixels within emission-line regions, camera edges (9 pixels
    each side), and pixels with zero inverse variance.
  \item Slide a window of width 75 pixels with step 125 pixels.  In each
    window, apply a $2\sigma$ iterative clip and compute the robust median
    and inter-quartile sigma of the unmasked flux values; discard windows
    with fewer than 15 good pixels.
  \item Fit a weighted quadratic spline (\texttt{UnivariateSpline}, $k=2$)
    to the window medians, using the per-window inverse robust sigma as
    weights.  The spline uses constant extrapolation (\texttt{ext=3}) beyond
    the wavelength range of the input windows.
\end{enumerate}

The resulting smooth continuum is added to the SED model before emission-line
fitting, so that the line fitter operates on residuals with the smooth
baseline already removed.


\section{Monte Carlo Uncertainty Estimation}
\label{sec:monte}

Parameter uncertainties are estimated by repeating the continuum fit on
$N_{\mathrm{monte}} = 10$ independent realizations of the data.  Each
realization is drawn independently from Gaussian distributions with means
equal to the observed fluxes (spectral and photometric) and standard
deviations equal to the measurement uncertainties:
\begin{equation}
  f_i^{(m)} \;\sim\; \mathcal{N}(f_i^{\mathrm{obs}},\,\sigma_i^2),
  \qquad m = 1,\ldots,N_{\mathrm{monte}}.
\end{equation}
The inverse variance of each fitted quantity $\theta$ is then estimated as
\begin{equation}
  \sigma_\theta^{-2} \;=\; \bigl[\mathrm{Var}(\theta^{(1)},\ldots,\theta^{(N_{\mathrm{monte}})})\bigr]^{-1}.
\end{equation}
Uncertainties are reported for $\tauv$, $\vdisp$ (when fitted), Dn(4000),
rest-frame luminosities, and observed-frame continuum fluxes.  The spectral
Monte Carlo realizations are drawn once (using a fixed random seed) and
reused for both the continuum and emission-line fits.


\section{Physical Output Quantities}
\label{sec:outputs}

\paragraph{Stellar mass.}
The total stellar mass is
\begin{equation}
  \Mstar \;=\; \mathcal{M}_0 \sum_k c_k\,m_{\star,k},
\end{equation}
where $m_{\star,k}$ is the surviving stellar mass fraction (in units of
$\Msun$ per $\Msun$ formed) for template $k$, as returned by FSPS.

\paragraph{Dn(4000).}
The 4000\,\AA\ break index Dn(4000) is computed from the line-free
best-fit model following Balogh et al.\ (1999):
\begin{equation}
  D_n(4000) \;=\;
  \frac{\int_{4000}^{4100}\!\flam^{\mathrm{model}}\,d\lambda}
       {\int_{3850}^{3950}\!\flam^{\mathrm{model}}\,d\lambda}.
\end{equation}
It is also measured directly from the observed spectrum (with inverse-variance
weighting) as a check on the model.

\paragraph{Rest-frame luminosities.}
Monochromatic rest-frame luminosities are measured from the line-free
best-fit SED at 1450, 1500, 1700, 2800, 3000, and 5100\,\AA:
\begin{equation}
  L_\nu(\lambda_0) \;=\;
  4\pi\DL^2\,(1+z)\,\flam^{\mathrm{model,nolines}}(\lambda_0(1+z))\,
  \frac{\lambda_0^2}{\C},
\end{equation}
measured by a local linear fit to the SED in a $\pm50$\,\AA\ window.
The values at 1500 and 2800\,\AA\ are expressed as $\nu L_\nu$ to facilitate
UV SFR estimates.

\paragraph{Continuum fluxes.}
Observed-frame continuum flux densities are measured adjacent to the
Ly$\alpha$, [O\textsc{ii}], H$\beta$, [O\textsc{iii}], and H$\alpha$ lines,
for use in equivalent-width computations during emission-line fitting.

\paragraph{K-corrections and absolute magnitudes.}
Observed-to-rest-frame K-corrections are computed by comparing synthesized
photometry from the best-fit SED at the observer-frame and rest-frame filter
effective wavelengths; absolute magnitudes follow from the distance modulus
and K-correction.

\paragraph{Reduced $\chi^2$.}
Separate reduced chi-squared values are reported for the spectroscopic and
photometric fits; the degrees of freedom account for the number of non-negative
template coefficients actually used (i.e., those with $c_k > 0$), plus
$\tauv$ (and $\vdisp$ when fitted).


\section*{Acknowledgments}

The VARPRO approach to separable least squares follows the classic framework
of Golub \& Pereyra (1973).  The FSPS code is due to Conroy, Wechsler \&
van Dokkum (2009) and Conroy \& Gunn (2010); the C3K stellar library and
MIST isochrones are described in Choi et al.\ (2016) and references therein.
The Draine \& Li (2007) dust emission templates, the Nenkova et al.\ (2008)
AGN torus model, the Vestergaard \& Wilkes (2001) Fe template, and the
Inoue et al.\ (2014) IGM model are used as described in the text.

\end{document}
